% По материалам KACTL; Anders Sjoqvist, Ulf Lundstrom.
\chapter{Математические формулы}
\section{Уравнения}
\[ ax^2+bx+c=0\Rightarrow x=\frac{-b\pm\sqrt{b^2-4ac}}{2a}. \]
Абсцисса экстремума: $x=-b/(2a)$.
\[\begin{aligned}ax+by=e\\cx+dy=f\end{aligned}
\Rightarrow x=\frac{ed-bf}{ad-bc},\quad y=\frac{af-ec}{ad-bc}.\]
Правило Крамера для $Ax=b$: $x_i=\det A_i/\det A$,
где столбец i в $A_i$ заменён на b, $\det A\ne0$.
\section{Рекурренты}
Если $a_n=c_1a_{n-1}+\dots+c_ka_{n-k}$ и корни
$x^k-c_1x^{k-1}-\dots-c_k$ различны, то $a_n=\sum_i d_ir_i^n$.
Корень r кратности m даёт слагаемое $P(n)r^n$, $\deg P<m$.
\section{Тригонометрия}
\begin{align*}
\sin(v\pm w)&=\sin v\cos w\pm\cos v\sin w,\\
\cos(v\pm w)&=\cos v\cos w\mp\sin v\sin w,\\
\tan(v\pm w)&=\frac{\tan v\pm\tan w}{1\mp\tan v\tan w},\\
\sin v+\sin w&=2\sin\frac{v+w}{2}\cos\frac{v-w}{2},\\
\cos v+\cos w&=2\cos\frac{v+w}{2}\cos\frac{v-w}{2}.
\end{align*}
Для противоположных сторонам V,W углов v,w:
$(V+W)\tan((v-w)/2)=(V-W)\tan((v+w)/2)$.
\begin{align*}
a\cos x+b\sin x&=r\cos(x-\phi),\\
a\sin x+b\cos x&=r\sin(x+\phi),
\end{align*}
где $r=\sqrt{a^2+b^2}$, $\phi=\operatorname{atan2}(b,a)$.
\section{Геометрические формулы}
\subsection{Треугольник}
Стороны a,b,c, противоположные углы $\alpha,\beta,\gamma$,
полупериметр $p=(a+b+c)/2$.
\begin{align*}
S&=\sqrt{p(p-a)(p-b)(p-c)},\\
R&=\frac{abc}{4S},\qquad r=\frac Sp,\\
m_a&=\tfrac12\sqrt{2b^2+2c^2-a^2},\\
l_a&=\sqrt{bc\left(1-\frac{a^2}{(b+c)^2}\right)},\\
\frac{\sin\alpha}{a}&=\frac{\sin\beta}{b}=\frac{\sin\gamma}{c}=\frac1{2R},\\
a^2&=b^2+c^2-2bc\cos\alpha,\\
\frac{a+b}{a-b}&=\frac{\tan((\alpha+\beta)/2)}{\tan((\alpha-\beta)/2)}.
\end{align*}
R,r — радиусы описанной и вписанной окружностей;
$m_a,l_a$ — медиана и биссектриса к стороне a.
\subsection{Четырёхугольник}
Стороны a,b,c,d, диагонали e,f с углом $\theta$;
$F=b^2+d^2-a^2-c^2$, площадь S:
\[4S=2ef\sin\theta=F\tan\theta=\sqrt{4e^2f^2-F^2}.\]
Для вписанного: противоположные углы в сумме $\pi$,
$ef=ac+bd$, $S=\sqrt{(p-a)(p-b)(p-c)(p-d)}$,
$p=(a+b+c+d)/2$.
\subsection{Сферические координаты}
$\theta$ — угол от оси z, $\phi$ — азимут от оси x:
\begin{align*}
x&=r\sin\theta\cos\phi,& r&=\sqrt{x^2+y^2+z^2},\\
y&=r\sin\theta\sin\phi,& \theta&=\arccos(z/r),\\
z&=r\cos\theta,&\phi&=\operatorname{atan2}(y,x).
\end{align*}
\section{Производные и интегралы}
\begin{align*}
(\arcsin x)'&=\frac1{\sqrt{1-x^2}},&
(\arccos x)'&=-\frac1{\sqrt{1-x^2}},\\
(\tan x)'&=1+\tan^2x,&(\arctan x)'&=\frac1{1+x^2},\\
\int\tan ax\,dx&=-\frac{\ln|\cos ax|}{a},\\
\int x\sin ax\,dx&=\frac{\sin ax-ax\cos ax}{a^2},\\
\int e^{-x^2}\,dx&=\frac{\sqrt\pi}{2}\operatorname{erf}(x),\\
\int xe^{ax}\,dx&=\frac{e^{ax}}{a^2}(ax-1).
\end{align*}
Интегрирование по частям, $F'=f$:
$\int_a^bfg=[Fg]_a^b-\int_a^bFg'$.
\section{Суммы и ряды}
\[\sum_{i=a}^bc^i=\frac{c^{b+1}-c^a}{c-1}\quad(c\ne1).\]
\begin{align*}
\sum_{i=1}^n i&=n(n+1)/2,\\
\sum_{i=1}^n i^2&=n(n+1)(2n+1)/6,\\
\sum_{i=1}^n i^3&=n^2(n+1)^2/4,\\
\sum_{i=1}^n i^4&=\frac{n(n+1)(2n+1)(3n^2+3n-1)}{30}.
\end{align*}
\begin{align*}
e^x&=\sum_{k\ge0}\frac{x^k}{k!},\\
\ln(1+x)&=x-x^2/2+x^3/3-\dots\quad(-1<x\le1),\\
\sqrt{1+x}&=1+x/2-x^2/8+x^3/16-5x^4/128+\dots,\\
\sin x&=x-x^3/3!+x^5/5!-\dots,\\
\cos x&=1-x^2/2!+x^4/4!-\dots.
\end{align*}
Ряд корня сходится при $|x|\le1$, ряды exp/sin/cos — для всех x.
\section{Вероятности}
$\mathbb EX=\sum_x xp(x)$,
$\operatorname{Var}X=\mathbb EX^2-(\mathbb EX)^2$,
$\sigma=\sqrt{\operatorname{Var}X}$.
Для непрерывной величины суммы заменяются интегралами с плотностью.
Всегда $\mathbb E(aX+bY)=a\mathbb EX+b\mathbb EY$;
для независимых X,Y:
$\operatorname{Var}(aX+bY)=a^2\operatorname{Var}X+b^2\operatorname{Var}Y$.
\subsection{Дискретные распределения}
\textbf{Биномиальное}: число успехов в n независимых испытаниях
с вероятностью успеха p.
$p_k=\binom nkp^k(1-p)^{n-k}$,
$\mu=np$, $\sigma^2=np(1-p)$.
При малом p приближается распределением Пуассона с $\lambda=np$.

\textbf{Геометрическое}: номер первого успеха, k>=1.
$p_k=p(1-p)^{k-1}$, $\mu=1/p$, $\sigma^2=(1-p)/p^2$.

\textbf{Пуассона}: число событий за время t при интенсивности
$\kappa$, $\lambda=\kappa t$.
$p_k=e^{-\lambda}\lambda^k/k!$, $\mu=\sigma^2=\lambda$.
\subsection{Непрерывные распределения}
\textbf{Равномерное на [a,b]}:
$f(x)=1/(b-a)$ внутри интервала и 0 вне;
$\mu=(a+b)/2$, $\sigma^2=(b-a)^2/12$.

\textbf{Экспоненциальное}: время до события Пуассона,
$f(x)=\lambda e^{-\lambda x}$ при x>=0, иначе 0;
$\mu=1/\lambda$, $\sigma^2=1/\lambda^2$.

\textbf{Нормальное} $\mathcal N(\mu,\sigma^2)$:
$f(x)=\exp(-(x-\mu)^2/(2\sigma^2))/(\sqrt{2\pi}\sigma)$.
Для независимых $X\sim\mathcal N(\mu_1,\sigma_1^2)$,
$Y\sim\mathcal N(\mu_2,\sigma_2^2)$:
\[aX+bY+c\sim\mathcal N(a\mu_1+b\mu_2+c,a^2\sigma_1^2+b^2\sigma_2^2).\]
\section{Цепи Маркова}
Следующее состояние зависит только от текущего.
$P_{ij}=\Pr(X_{t+1}=i\mid X_t=j)$, $p^{(n)}=P^np^{(0)}$.
Стационарное распределение удовлетворяет $\pi=P\pi$.
Для неприводимой цепи $\pi_i=1/\mathbb ET_i$, где $T_i$ — время
возврата в i. Отношение $\pi_j/\pi_i$ — ожидаемое число посещений j
между возвратами в i (с соответствующим учётом граничного посещения).
В связном неориентированном графе случайное блуждание по соседям
имеет стационарные вероятности, пропорциональные степеням.
Для сходимости из любого старта нужны неприводимость и
апериодичность (НОД длин возвратов равен 1);
тогда $P^n\to\pi\mathbf1^T$. У связного неориентированного
недвудольного графа эти условия выполнены.

Для поглощающего состояния i вероятность попасть в него из j:
$a_{ij}=P_{ij}+\sum_{k\in G}a_{ik}P_{kj}$,
где G — непоглощающие состояния.
Среднее время до поглощения из i:
$t_i=1+\sum_{k\in G}P_{ki}t_k$.
Предполагается, что из каждого состояния G достижимо поглощение.

\section{Частичный порядок: цепи и антицепи}
Отношение $\le$ — частичный порядок, если оно рефлексивно,
антисимметрично и транзитивно. Элементы сравнимы, если $a\le b$
или $b\le a$.
\textbf{Цепь}: подмножество, в котором любые два элемента сравнимы.
\textbf{Антицепь}: любые два различных элемента несравнимы.
\textbf{Теорема Дилуорса}: в конечном частично упорядоченном множестве
максимальный размер антицепи равен минимальному числу цепей,
на которые можно разбить всё множество.
\textbf{Двойственное утверждение (Мирский)}: максимальная длина цепи
равна минимальному числу антицепей в разбиении.

Как посчитать: две копии всех n элементов, ребро $u_Lv_R$ при $u<v$.
Если максимальное паросочетание имеет размер M, ответ равен $n-M$.
Совпавшие пары задают ссылки на следующий элемент цепи; начинаем
с элементов без предшественника.
Для порядка по достижимости в DAG нужны рёбра \emph{транзитивного
замыкания}, а не только исходные рёбра! Паросочетание по исходным
рёбрам даёт покрытие непересекающимися ориентированными путями;
это другая задача. Пример: при порядке делимости
$\{1,2,3,6\}$ цепь $1<2<6$, антицепь $\{2,3\}$, ответ 2.
Эти цепи не связаны с вероятностными цепями Маркова.

\section{Римские числа}
\kactlimport{RomanNumerals.h}

\section{Расстояние по сфере}
\begingroup
\kactlcontentdir{content/geometry}
\kactlimport{sphericalDistance.h}
\endgroup
